\documentclass{ltjsbook}
\usepackage{docmute} 
\usepackage{../../myStyle} 

\begin{document} 
\chapter{モデル比較}
\label{x30_comparison} 
さて，今回でこの授業も最終回となりました。今回は授業全体のまとめとして，これまでの流れを俯瞰的に捉え直すとともに，今後みなさんがこの講義を通じて学んだことが，どのように利用できるのかに言及していきたいと思います。

本書の前半，心理学データ解析応用1として前期にお話ししてきた内容は，心理学で利用されるデータ，特に心理尺度を用いたものが，どういう根拠で「心」を測定していると言えるのかについて解説しててきました。
後半，心理学データ解析応用2の方は，プログラミングやベイズ統計といった目新しい研究手法の方が目についたかもしれませんが，その背後にあったメッセージは，どういうメカニズムでデータが生成されてきたのかという，\textbf{データ生成メカニズム}というアイデアをもつことでした。その意味で，本書全体のメッセージは一貫しています。すなわち，心理学者がデータを取るという時の，データに込められた意味や想定されるメカニズムをどのように表現し抽出し解釈するかという，心理学的営みそのものをお伝えしたかったのです。ここに無知・無自覚なまま，心理尺度で何かが測定できると考えたり，検定の結果から意味を解釈したりできるはずがないのですから。

そのための準備として，確率統計に関する理論や，コンピュータ，プログラミングに関する技能などの習得が必要だったわけです。心理学を学ぼうと思って入ったはずの大学で，どうして数学やプログラミングをしているのか，こんなはずじゃなかったと思った人もいるかもしれませんが，逆にどうして既知の知識や技術だけで，この摩訶不思議な心，人間，社会のありようが理解できると思っているのかと言いたいほどです。もちろん他のアプローチや，他にも学ぶべきことがたくさんありますし，時間は有限ですから苦手な手法は後回しにしたくなる気持ちはわかりますが，実は他の教養を身につけるよりもこちらの方が直接的で近道なルートだといえるかもしれません。

\section{ベイジアンモデリング}
Stanを使った統計的アプローチは，\keyterm{ベイジアンモデリング}{ベイジアンモデリング}{Bayesian Modeling}と呼ばれることがあります。Stanが事後分布からの乱数発生機であり，事前分布と尤度を記述するだけで事後分布を得ることができるというところが，ベイズの定理に根拠を持ったベイズ統計学ですから，「ベイジアン」という冠がつくのは当然です。しかし後半の「モデリング」については，なにもベイズ統計学に限った話ではありません。モデルを立てて考えるということだけを考えれば，広い意味では心理学だって「心」というモデルを立てて考えているのですから。もっというと，学問というのは一般的に抽象化(あるいは単純化)した理想の世界での論理的構造を扱うものですから，全てモデリングアプローチだと言えるかもしれません。

心理学はその研究方法として，調査，実験，観察，介入といったものがあります。これらを通じて心とはなにかということを考えていくわけです。特に実験や介入においては，条件の制御や統制を行って，それに伴って生じた結果の違いを\textbf{効果}と考えるのでした。そのことと\keytermL{帰無仮説検定}{きむかせつけんてい}は非常に相性が良かった事は，想像に難くないでしょう。帰無仮説検定は，出てきたデータと変数の関係について，機械的に判断を下すことができる方法論でした。帰無仮説検定が農学から生まれてきたことからもわかるように，その方法はどの研究分野に対しても適用できる，共通のルール・御作法なのです。心理学者は\underline{実験計画に心理学的要素を埋め込み}，結果を統計的に判断することで，心理学的な結果の解釈を進めるてきました。いいかえれば，実験の計画を立てることそのものが心理学のエッセンスであったのです。

しかし皆さんはすでに，実験計画が統計的に見れば\keytermL{線形モデル}{せんけいもでる}にすぎないことを知っていますね。そしてごく単純な線形モデルによってしか，心理学のエッセンスを表現する方法がなかった時代にあっては，想定しうる心のメカニズムが平均値の差として出てくるように置き換えるほかなかったのです。方法論に自由がない分，その他のところで工夫するしかなかったとも言えます。そうした工夫の中には芸術的なまでに作り込まれているものがあったりしますから，そこに痺れる憧れる，という人も少なからずいるのは当然ともいえるでしょう。

線形計画の結果は，操作/介入の効果があったかなかったかという判断になりがちです。結果がそれしか示しておらず，その結果を生むための工夫はすでになされており，得られるデータの精度も平均値差を示すことができれば良い程度にしかない，ともいえるかもしれません。しかし測定のツールは時事刻々と進歩してきましたので，より精緻に違いを見て，よりリッチな解釈を許す意味合い豊かなデータも得られるようになってきました。そのような時，伝統的な方法論だけで立ち向かうのは勿体無いのではないでしょうか。
方法論の呪縛から自由になり，さまざまな表現ができるようになれば，心理学ももっと発展させることができるでしょう。

モデリングは，変数間関係を数式で表現することでもあります。「心」や「気持ち」といった表現を始め，「幸せ」とか「恋愛感情」といった日常用語で考えるべきところを，$X_i,Y_{ij},\alpha,\beta...$といった記号で置き換えて表現するのには慣れが必要です。ひょっとしたら抵抗を感じる人もいるかもしれません。しかし，数式で表現する事は，他の解釈を許さない一意的表現をすることでもあります。人間って色々だよね，当てはまることもあれば当てはまらないこともあるよね，という日常的な理解で留まりたいのであれば--そしてそこで止まっていられるのは幸せなことでもあるのですが--それでいいのですが，学問を進める以上はより明確に，より誤解のない表現をしなければなりません。個別のケースにしか当てはまらないこと，その程度を確率で表現して，理論的に全体的な傾向を「こうなっているに違いない」と厳密に表現するには，日常用語ではなく数式用語が必要なのです。モデリングを進める上で，そうした表現ができるようにトレーニングする必要はありますが，使えるようになると折線回帰や状態空間モデルのように，さまざまな表現ができることも見てきた通りです。

この授業を受けてきたからと言って，さあ今すぐ自分で心理学的モデルを作れと言われても難しい，と思うかもしれません。実際これまで私が教えてきた中で，よく質問されるのは「やり方はわかるけれども，自分でやれる気がしない。どうやってモデルを思いついたらいいのか」ということです。これに対する答え方は2つあって，1つは「色々なパターンを見て模倣する」です。「学ぶ」は「真似ぶ」からきているとも言われたりしますが，どのような表現の可能性があるのかについては，実際の応用パターンをさまざま見るのが早いでしょう。馴染みのない料理や調理法があれば，レシピ本を買って応用パターンをいろいろ仕入れますよね。そうすることで，「こんな工夫ができるのか。じゃあ自分はこうしてみよう，ここをちょっと変えてみよう」と前に進むことができるのですから。こうしたモデルがたくさん載っているものとして，\citep{Lee201709}や\citep{豊田秀樹2017,豊田秀樹2018,豊田秀樹2019}などがありますので，パラパラっとみながら応用例を広く眺めてみるといいでしょう。

もう1つの答えは，「紙とペンをつかって設計図を書く」というものです。この方法は，この講義の中で再三お伝えしてきたやり方になります。データを見た時に，どういうメカニズムでデータが得られたかを考えるところにこそ，心理学者の本当の興味関心があるはずです。このデータは心のこういう仕組みを反映しているのではないか。データの背後には変数同士のこういう関係が潜んでいるのではないか，ということを考えていくわけです。わからないところは確率で表現しつつ，変数間関係を関数で記述する。そのための設計図です。設計図とStanがあれば答え(事後分布)は得られるのですから，データ生成の仕組みを作り込むことに我々は全力を傾けることができます。複雑な関数関係をいきなり思いつくのが難しい，ということもあるでしょう。その場合はデータを可視化し，データに合うような関数はどんな形なのかを考える，あるいは一旦，単純な線形モデルを当てはめて，それを徐々に複雑にしていく，というやり方が良いでしょう。\citet{Maeda2017}は心理学者が書いたベイズ統計の本ですが，その後半は線形モデルです。一般化線形モデルや階層モデルについて，特に複雑な前置きなく「こういう仕組みになってるんだから，とりあえずそれを反映して，関係は線形でも大体うまくいくでしょ」という感じで話が進んでいきます。簡単なモデルから徐々に複雑にしていく，というのがモデリングの王道でもあります。完成品だけをみるととても複雑に思えますが，設計図片手に紐解いていけば，意外とわかりやすいかもしれません。

モデリングは自由に設計図を書くことができるので，分散分析やt検定など平均値差だけにこだわっている方法をみると，そちらがいかにも窮屈そうにしていると思えるかもしれません。講義を通じて得た知識や技術で，自由にデータを調理していただければと思います。

\section{帰無仮説検定の代案}
ベイジアンモデリングは，自由に書いた設計図から確率的プログラミング言語をつかって答えを得る方法です。プログラミング言語が使えるようになるの大変です。これまでの一般的な統計ソフトでは，ボタン1つで答えを出してくれたのに，と思う人もいるかもしれません。しかしそれは，従来の分析方法がさまざまな仮定や前提に基づいた，マニュアル的アプローチを許す範囲のものに限定されていたからですし，「簡単，楽ちん」ということが「知らなくていい」ということだとすると，その考え方がさまざまな問題を引き起こしてきたことを知らなければなりません。これは\textbf{心理学における再現性問題}\citep{Ikeda2016}として取り沙汰されており，心理学そのものに疑問符を投げかけるような事態を招いています。その原因の1つが，統計的技術の誤用・悪用にありました。

心理学ではながらく，帰無仮説検定によって理論を積み重ねてきました。帰無仮説検定では$p$値が5\%より小さければ良し，という「ここだけ見ておけば良い」というようなマニュアル的基準があり，差があればすごい効果が発見されたぞ！と喧伝するということで進んできた側面があります。すでに述べたように，この\keytermL{有意差}{ゆういさ}を出すために要因計画の方にこそ心理学者は注力してきたのですが，$p<0.05$になりさえすれば良いという表面的な理解しかしていないマニュアル人間が出来上がってしまうことが，問題を大きくしたのかもしれません。

ベイズ統計によるアプローチの場合，データがどのように出てきたのか，事前分布はどうなのか，と言ったことを考えないわけにはいきません。自由に設計図を描くことができるというのは，裏を返せば，最初は白紙でしかないということであり，きちんと動く統計ツールを作るためには細部まで全て記述しなければならないのです。帰無仮説検定では，極端な場合「データ生成には正規分布が仮定されている」ということを知らなくても，$p$値がどうなるかだけをみる事はできてしまうのに，です。このように，ベイズ統計を使う利点の1つは，前提や仮定について全て自覚的に，明示的に記述しなければならないというところにあり，これが悪用・誤用を生みにくくする安全装置となっている面があります。

最近ではJASP\citep{JASP2021}という，GUIでベイズ分析もできるソフトウェアが出てきました。帰無仮説検定と同じような分析を，ベイズでやったらどうなるか，すぐに試すことができます。この便利なアプリではStanやJAGSなどとは違って，どのような前提があるかを自覚しなくても，分析結果を得る事はできます。ですから，今挙げた「安全装置」としてのベイズ統計の利点は，簡単に外れてしまう日が来るかもしれませんし，その方がユーザにとってはいいかもしれません。
この講義では第\ref{x19_modeling1}講から\ref{x22_modeling4}まで，帰無仮説検定の代わりにモデリングアプローチをする例を紹介してきました。
さて，ではこうした帰無仮説検定の代わりとして，ベイズ統計を使う利点はどこにあるでしょうか。

もちろんいくつか考えられますので，箇条書き的に解説してみたいと思います。

\paragraph{幾重にも組み合わさる仮定や補正からは無縁} たとえば二群の平均値の差を検定する場合，データが正規分布に従っているかどうかを検定し，また二群の分散が等しいかどうかも検定して，問題なければt検定に進みます。もし分散が等しくないということになれば，Welchの補正をしなければならない，とマニュアルにはあります。またたとえば，WithinのANOVAをするときも，球面性の検定を行なってデータが分析の前提となる仮定にあっているかどうかを判定しなければなりません。主効果が見られた，交互作用が見られたとなれば，今度はどこに差があるのかをみるために--検定の繰り返しを避けるために，下位検定を行います。

このように，仮説検定をする場合は条件にあっているかどうかを検定し，あっていなければ補正をするなどの手続きが必要です。同時に，仮説検定を同一のデータに何度も繰り返す事は，5\%水準というタイプ1エラーを生じる確率が適切に制御できなくなるため，本来なら実験単位ではなく「一連の分析単位」でこれを調整しなければなりません。正直なところ，ここまで厳密な調整が，あらゆる心理学研究できちんと行われているかと言われれば，その答えはNoなのですが。

ベイズ統計の場合，こうした問題は非常に軽微なものに変わります。たとえば二群の分散が等しいかどうかについては，「等しくないモデル」を考えれば良いのです。具体的にいうと，二群の分散が等しいモデルは，
$$ X_i \sim N(\mu_1,\sigma), Y_i \sim N(\mu_2,\sigma)$$
であり，等しくないモデルは
$$ X_i \sim N(\mu_1,\sigma_1), Y_i \sim N(\mu_2,\sigma_2)$$
とするだけです($\sigma$に添字がある，すなわち違う大きさとして推定する)。球面性の検定についても分散共分散行列の要素それぞれを推定するだけであり，その数値に事前に定められた制約はありません。

またたとえば，タイプ1エラーの制御についても，そもそもタイプ1エラーという発想がありませんから存在しないのです。データがあって，平均構造を線形モデルで表現したのであれば，そこから出てくる事後分布は1つしかありません。その単一のパラメータ同時事後分布をどのように切り分けようとも，ある判断のエラーに確率が伴う事はないのです。

帰無仮説検定を厳密に行うには，下位検定を事前に計画してどこにどのような差があると考えるかも準備しておく必要があります。帰無仮説検定は勝敗を決める手法ですから，ゲームのルールは事前に決めなければならないのです。ゲームの設定如何によっては，判定結果が変わってくることがあるからです。これに対してベイズ統計的アプローチであれば，結果(事後分布)は決まっているので，下位検定を事前に計画する必要はなく，事後分布をどう切り分けようと分析者の自由です。
\paragraph{停止規則やサンプルサイズの問題が生じない}
同様のことは，サンプルサイズの設計についても言えます。帰無仮説検定の場合は，ゲームのルールを事前に決める必要があると言いました。このルールの中には，サンプルサイズも含まれています。サンプルサイズが大きくなればなるほど，有意であると判定される可能性は上がっていきます。帰無仮説検定における，最もやってはいけない研究実践のひとつは，データを取りながら検定を繰り返し，有意になったらデータを取るのをやめる，というものです。検定を繰り返すと，誤った判断をしてしまう可能性が増えますし，サンプルサイズが大きくなると有意になりやすくなります。つまり「勝つまでゲームをやめない」というのは卑怯な方法なのです。

しかし実際は，数人相手にデータをとってみたけど有意差はでなかった，でも惜しいから頑張ってデータを増やした，結果としてある日有意な差になったので，サンプルサイズをちゃんと記載して論文として提出した，ということがあるわけです。努力して真実に辿り着くのはいい話ですが，この例は努力したことが結果の捏造になっていたかもしれないという，より悲しい話になってしまいます。

帰無仮説検定は事前にサンプルサイズを決めなければいけません。このサイズ，この基準で勝負するぞ，というルールがあってからの一発勝負なのです。どこでデータの収集を止めるかというのは，\keyterm{停止規則}{ていしきそく}{Stopping Rule}といって，これまた事前に決めておかねばならないことなのです。事後的に，勝負に勝ったからこのやり方でよかったんだろう，と考えるのは間違っているのです。

これに対して，ベイズモデルの場合はこうした問題が生じません。最後が確率的な判断にならないので，こうした問題が生じないのです。データが蓄積されていくことは，徐々に確信度が高まっていくことでもあります。差があるのかないのか，という判断をするにあたって，よりその違いが明確になっていくだけなのです。

\paragraph{より柔軟な仮説を立てることができる}
帰無仮説検定は，帰無仮説と対立仮説という2つの仮説の比較判断になります。ここで帰無仮説は「差がない」「相関がない」といったものになります。たとえば二群間の母平均を比較する場合，帰無仮説は$\mu_1 = \mu_2$になります。
しかし実際には，差がないということを主張したい場合もあるかもしれません。帰無仮説に「差がある」というのをおくことができないのは，「あるというのは，どれぐらいあればいいのか」という問題がすぐに浮上するからです。たとえば母平均の差が$\mu_1 - \mu_2 = 0.1$なのか，$\mu_1 - \mu_2 = 0.11$なのか，$\mu_1 - \mu_2 = 0.111$なのか・・・言い出すとキリがないですね。$p$値は帰無仮説のもとで計算されるものですから，差があるという仮説は無限に作れてしまうので，検証すべき帰無仮説も無限に膨れ上がってしまいます。差がないという状態は$\mu_1 - \mu_2 = 0$と一意に定まるので，帰無仮説を設定することが可能なのです。

では帰無仮説を積極的に支持すればよいではないか，と思われるかもしれませんが，それができません。判定結果の$p$値は帰無仮説のもとで計算していますから，「帰無仮説のもとで考えるとおかしな結果になった」という背理法が成り立たないのです。「帰無仮説のもとでおかしな結果にならなかった」というのは，帰無仮説以外の可能性を除外できるものではありません。

このように，帰無仮説検定では「効果がなかった」とか「関係がなかった」とは言えず，せいぜい「効果がないとは言えない」というにとどまるのでした。
これに対して，ベイズ統計では\keytermL{モデル比較}{もでるひかく}の観点から検証することになります。帰無仮説は$\mu_1=\mu_2$というモデル，対立仮説は$\mu_1 \neq \mu_2$というモデルです。言い換えれば，帰無仮説は
$$ X_i \sim N(\mu_1,\sigma_1), Y_i \sim N(\mu_1,\sigma_2)$$
というモデルであり，対立仮説は
$$ X_i \sim N(\mu_1,\sigma_1), Y_i \sim N(\mu_2,\sigma_2)$$
というモデルだとしているに過ぎないからです。

これまでみてきたように，ベイズ統計のアプローチで平均値の差を考える場合は，$\mu_1 - \mu_2$のような差の分布を\keytermL{生成量}{せいせいりょう}でつくり，その大きさをそのまま吟味できます。どうしても「差があるのか，ないのか」という判断がしたければ，\keyterm{実質的に等価な範囲}{じっしつてきにとうかなはんい}{Region Of Practical. Equivalence; ROPE}($\to $ セクション\ref{ROPE}, Pp.\pageref{ROPE})を事前に設定し，その区間に差の分布が入ってくるかどうかを判断するのでした。そもそも差も分布しているわけですから，どれぐらい重複しているかという幅で考えるほかないわけです。帰無仮説検定の文脈においても，\keyterm{効果量}{こうかりょう}{Effect Size}とよばれる\textbf{標準化された差}の大きさを算出して考えますが，このように「どれぐらい違いがあれば差があると言えるか」という程度問題が前面に出てくることが重要なのです。5\%水準のように機械的に判断するのではなく，領域固有の知識(ドメイン知識)に基づいて，実質的な差の大きさをケースバイケースで判断することが重要です。もちろん機械的に手続きを進められないことは不便かもしれませんが，不便だからといって真実を曲げるようでは本末転倒です。

また生成量を使ったアプローチのところで解説したように，パラメータについて考えるだけでなく，そのモデルから生成されるであろうデータについて仮説を立てたり検証したりできるのも，ベイズ統計的アプローチの利点です。実際どの程度の差があると意味があるのかについて，具体的な数値シミュレーションで考えることができるからです。パラメータが$\delta$以上の差があれば良いとか，二群のデータが$D$以上違ってくる確率はどれぐらいか，といった検証の仕方は，仮想空間上の$p$値よりも具体的で意味のある情報をもっています。パラメータやデータについて，「同じかどうか」以上の情報を検証できることは，ベイズ統計の長所といえるでしょう。

\section{モデル比較}
先ほど，\textbf{モデル比較}ということに言及しました。このことについて，最後に少し触れておきたいと思います。
ベイジアンモデリング，すなわちモデルをたててベイズ統計学的に推定するなかで，そのモデルの優劣を考える方法があります。

もう一度ベイズの公式を見てみましょう。ベイズの公式は次のようなものでした。

\[ p(\theta| D) = \frac{p(D| \theta)p(\theta)}{p(D)} \]

ここで右辺の分母は$p(D)$，つまりデータの確率をあらわすもので\keyterm{周辺尤度}{しゅうへんゆうど}{marginal likelihood}と呼ばれます。
これはどのように計算されるのでしょうか。尤度は$p(D| \theta)$で表されますが，$P(D)$はこのパラメータが無くなったものになります。どのようにしてパラメータをなくすかというと，パラメータのありそうな可能性を全て足し合わせることで\keytermL{周辺化}{しゅうへんか}すれば良い，ということになります。

たとえば性別と学年のクロス表があったとして，ランダムに選んだ1人の学生が女性である確率は，全ての学年の女性の確率を足し合わせたものになるわけです。記号で考えると，条件付き確率$P(D| \theta)$において，$\theta$の可能性を全部考えれば$P(D)$だけになるということです。これはパラメータ$\theta$が離散的な例ですが，連続的な場合は
$$ p(D) = \int p(D | \theta ) d\theta $$
のように積分で考えればよいでしょう。

これはモデルに含まれるあらゆるパラメータのパターンを網羅して，データが得られる確率を計算したことになりますから，モデルが表現できる世界の中でどれぐらいデータを捕まえることができたか，ということを表しているとも言えます。
ここでモデルが複数ある場合を考えましょう。あるモデルを$\mathcal{M}$と表現することにして，第一のモデルを$\mathcal{M}_1$,第二のモデルを$\mathcal{M}_2$，と表したとします。
すると事後分布も「あるモデルを仮定した時の事後分布」ということになりますから，次のようにすべて「モデルのもとでの」という条件付き確率で表現できます。

$$ p(\theta | D,\mathcal{M}_i) = \frac{p(D | \theta,\mathcal{M}_i)p(\theta | \mathcal{M}_i)}{p(D| \mathcal{M}_i)} $$

ここで右辺の分母を見ると，$p(D| \mathcal{M}_i)$，つまり「モデル$i$のもとでデータが得られる確率」を表していることになります。もしこの確率が高ければ，このモデルはこのデータを生む確率が高いわけです。逆にもしこの確率が低ければ，このモデルはこのデータを生む確率が低いということになります。この周辺尤度は，モデルがデータを支持する度合いであり，証拠(\textbf{エビデンス})の強さとも言えるわけです\footnote{これを考えると，どんなデータでも作れる幅広いモデルを考えておけばいいじゃないか，と思うかもしれませんが，モデルを広く考えすぎるとそこから出てくるデータはそのごく一部でしかなく，かえってエビデンスとしては弱くなってしまいます。なんでも予測できるモデルは何も予測していないことと同じなのです。}。

その上で，次のような数字を考えます。

\[ BF_{12} = \frac{p(D| \mathcal{M}_1)}{p(D| \mathcal{M}_2)} \]

この数字は\keyterm{ベイズファクター}{べいずふぁくたー}{Bayes Factor}と呼ばれるもので，モデル2に比べモデル1がどれほどデータに支持されているかを示す数字です。この数字が1.0より大きければ，モデル1はモデル2に比べて相対的に強くデータに支持されているということになります。逆に1.0を下回ると，モデル2のほうがモデル1よりもいいね，ということです。モデル同士の比による表現ですので，$BF_{21}$としても構いません。ここでモデル1が対立仮説，モデル2が帰無仮説だとすると，どちらがよりデータをうまく表しているか，どちらがよりデータに支持されているかがわかるわけです。最近では，帰無仮説検定の文脈でもBFを併記して，どちらのモデルが良いかを示すこともありますし，この方法ですと帰無仮説(というか差がないというモデル)を積極的に支持する，ということもできます。

この手法はベイズ推定するものであれば一般的に通用する考え方ですので，1つのデータに対していろいろなモデルを考えついたとしても，ベイズファクターを計算して相互に比較できるわけです。すなわち，ベイズのモデル比較はベイズファクターという共通の基盤を手にしたとも言えます。
そんな優れた方法があるのであれば，なぜもっと早く教えてくれなかったのか，という声が聞こえてきそうです。実は理屈上ではこの通りなのですが，この方法には実践上の問題が残されているのです。それは，\underline{周辺尤度の計算が難しい}というものです。全てのパラメータについて，考えられる全ての可能性を計算して足し上げるわけですから，それはそれは大変な計算になります。有効な近似計算の方法として，全パラメータの全確率空間を計算し尽くすかわりに，確率が高そうなところを効率よくサンプリングして推定値を得る\keyterm{ブリッジ・サンプリング}{ブリッジ・サンプリング}{Bridge Sampling}という手法が考えられており，その計算をするRのパッケージ\citep{bridgesampling}が提供されていますから，今後こうしたモデル比較も増えてくるでしょう。

またモデル評価の方法として，\keytermL{情報量規準}{じょうほうりょうきじゅん}によるアプローチも考えられています。これは「良いモデルとは良い予測をするモデルである」と考えるところから始まっており，平均的に良い予測をするモデル程度を指標化するものです。ベイズの定理を基本としながら，データを真のモデルから生成されるサンプルと考えたり，確率分布同士の近さを\keyterm{カルバック・ライブラー情報量}{かるばっく・らいぶらーじょうほうりょう}{Kullback-Leibler divergence}で表現するなど，非常に高度な数学的体系が確立されています。詳しくは\citet{浜田201912}など丁寧な解説書があり，Rパッケージなどで簡単に計算できるようになっています\citep{loo}。数学的にやや複雑な内容を含んでいますので，本講義の中では扱いきれんませんが，より進んだ研究や利用のために言葉の紹介だけはしておきます。


\section{おわりに}
本講ではここまで，ベイズ統計と従来の統計的分析を比較して議論してきました。ベイズ統計の方がメリットが大きいような解説をしてきましたが，これは筆者の好みもかなり含まれるものです。帰無仮説検定はなにも悪いものではなく，前提や仮定をしっかり守れば，一般的なソフトウェアで誰にでも結論を出すことができるという点で，とても良いものなのです。前提や仮定が多くて面倒だから，マニュアル人間ができてしまうというのは，あくまでもユーザ側の問題であって方法論の問題ではありません。ベイズ統計だって，モデルが間違っていれば結論は間違ったものになりますし，なによりデータに合わせてモデルを逐一設計図から書き起こすのは，大変コストがかかることでもあるからです。本来，これら2つの流儀は相反するものではなく，それぞれ依って立つ理論や考え方があって，必要に応じて選べるようになるのが一番です。

またベイズ流の研究アプローチ，特にモデル比較については，まだまだ議論が活発に行われている段階であり，特に心理学研究にこうしたアプローチがどれほど有用かについては，いまだに結論が出ていないところでもあります。
パラメータの推定値で考えるのが良いのか，モデル比較をして議論を進めるのかについても，まだまだ定番の方法というのは見つかっていません。

そもそも心理学の知見を数式に落としてモデリングできるようにするためには，心理学そのものの理論的発展や，人間についてのより深い理解も必要です。
人間の行動と予測が完全にできるようなモデルや一貫した理論体系は，まだまだ遥か先の目標に過ぎないというのが現状でしょう。
それでも有用なツールを武器に，少しずつ知見を積み重ねていくことができれば，心理学をより楽しく学べるのではないかとおもいます。Enjoy!



\end{document} 
